In this paper, we assume that the attacker targets a NUCA architecture with a multi-hop mesh interconnect.

\textbf{Attacker Goal:}
The attacker aims to infer secret-dependent memory access patterns of a victim program by observing LLC access latency differences in a NUCA architecture.

\textbf{Attacker Capabilities:}
The attacker can identify a vulnerable access pattern either through binary profiling or directly from the source code if available.
The attacker can time sensitive operations in the victim program through interaction with the victim.
Examples include sending or receiving messages through the network or IPC, detecting modifications of shared variables, or using other side channels.
The attacker may or may not have access to an accurate timing source (e.g., \texttt{rdtsc}); if not, alternatives such as counting threads can be used instead.
We also assume that the attacker does not have root privileges.

\textbf{Attacker Knowledge:}
The attacker knows the general microarchitectural properties of the target system, including the NoC topology and LLC bank mapping, but need not know the source code of the victim.

\textbf{Trust Assumptions:}
The proposed defense techniques assume a trusted CPU and OS.
Defending against a malicious CPU or OS is outside the scope of this work, as it would require a fundamentally different threat model (e.g., trusted execution environments).

\textbf{Scope and Limitations:}
This work does not defend against side channels in which the non-uniformity of the cache architecture is not the source of the timing difference.
The following are considered out of scope:
\begin{compactitem}
    \item Timing channels based on cache hit/miss differences~\cite{osvik2006cache};
    \item Inference channels based on cache eviction~\cite{yarom2014flush+,irazoqui2015s,liu2015last,gruss2015cache} and replacement policies~\cite{lru-sidechannel};
    \item Timing channels based on data-independent execution time differences in single-core architectures.
\end{compactitem}
For these side channels, we assume the victim program is already protected using existing solutions, such as cache partitioning~\cite{kiriansky2018dawg} or oblivious execution~\cite{oram,phantom}.

Finally, a transient execution attack~\cite{spectre,meltdown} could potentially exploit this side channel to leak secrets obtained during transient execution.
However, such an attack faces a fundamental challenge: the transient execution window is short and CPU microarchitectural noise makes timing unreliable.
We therefore consider transient execution attacks out of scope and leave their interaction with NUCA timing channels as future work.
